\section{The coefficients are exponents of the original divisor}
\label{sec:hironaka-snc-component-coefficients}

A sequential extraction produces genuine nonnegative coefficients, but
their meaning as exponents of the original divisor needs a separate
proof.  Here an actual finite product of powers of distinct SNC members
is first put in a genuine coordinate normal form.  Its original
divisor exponent on each member is then computed.  This calculation
identifies the coefficients produced by extraction, including zero
coefficients and points where several members meet.

\begin{lemma}[A power of a genuinely principal stalk]
\label{hir:component-power-principal-stalk}
\uses{hir:monomial-factorization}
\lean{Hironaka.HironakaSNCComponentCoefficients.idealPow_mem_iff_of_generator}
\leanok
Suppose an ideal $I$ is generated at $x$ by a holomorphic germ $u$.
For every integer $a\ge0$, the actual stalk $I_x^a$ is generated by $u^a$.
\end{lemma}
\begin{proof}
For $a=0$, both are the holomorphic unit ideal generated by one.
Inductively $I^{a+1}=I^a I$, and products of two genuinely principal
stalks are generated by the product of their generators.
Indeed every generator of the product ideal is a product of a
holomorphic multiple of $u^a$ with a holomorphic multiple of $u$;
finite holomorphic linear combinations remain multiples of $u^{a+1}$.
Conversely $u^a$ and $u$ belong to their actual factors, so $u^{a+1}$
belongs to the product and all its holomorphic multiples do as well.
\end{proof}

\begin{lemma}[The genuine finite product of principal stalks]
\label{hir:component-product-principal-stalk}
\uses{hir:component-power-principal-stalk,hir:divisor-finite-component-product}
\lean{Hironaka.HironakaSNCComponentCoefficients.componentPowerProduct_mem_iff_of_generators}
\leanok
Let $A=((W_i,a_i))$ be a finite list, and suppose each actual power
$(I_{W_i}^{a_i})_x$ is generated by a holomorphic germ $v_i$.
Then the complete stalk of the actual product ideal is generated by
$\prod_i v_i$.
\end{lemma}
\begin{proof}
Induct on the ordered list $A$.  The empty product is the holomorphic
unit ideal, generated by one.  For a nonempty list, the first power
ideal is generated by $v_1$, and the induction hypothesis gives the
generator $\prod_{i>1}v_i$ of the tail product.
The genuine principal-stalk product calculation from the preceding
proof supplies the generator $v_1\prod_{i>1}v_i$ of the full product.
The equality is membership in the complete actual product stalk,
including finite holomorphic combinations.
\end{proof}

\begin{lemma}[The actual weighted SNC product has a coordinate generator]
\label{hir:component-product-coordinate-generator}
\uses{hir:component-product-principal-stalk,hir:snc,hir:monomial-case,hir:hadamard-bridge}
\lean{Hironaka.HironakaSNCComponentCoefficients.exists_generator_componentPowerProduct}
\leanok
In positive ambient dimension, let $A=((W_i,a_i))$ be a finite weighted list whose underlying members
are distinct and SNC.  At every original point $x$, there are centred
adapted coordinates $\psi$ and distinct coordinate indices $\sigma(i)$
for the members through $x$, such that near $x$ the actual ideal
$\prod_i I_{W_i}^{a_i}$ is generated by
\[
 m=\prod_{i:x\in W_i}(\psi_{\sigma(i)})^{a_i}.
\]
\end{lemma}
\begin{proof}
Choose an actual adapted SNC chart at $x$.  Closedness and finiteness
permit an open shrinking avoiding every member not through $x$.
For each member through $x$, Hadamard's lemma identifies its actual
vanishing ideal with the principal ideal of its corresponding
coordinate function $u_i=\psi_{\sigma(i)}$ throughout the chart.
Lemma~\ref{hir:component-power-principal-stalk} identifies its actual
$a_i$th power with $(u_i^{a_i})$.
For a member not through $x$, its vanishing ideal is the unit ideal
on the chosen shrinking: every holomorphic germ vacuously vanishes
on that member near the point.  One generates all its powers.
The preceding product lemma gives a genuine generator of the entire
product.  Factors for inactive members are one and disappear, leaving
the displayed coordinate monomial.
Distinctness ensures that the coordinate assigned to a list member
is unambiguous; SNC gives distinct coordinates for the active members.
This produces the full nearby stalk presentation from the actual
vanishing ideals, rather than assuming a monomial presentation.
\end{proof}

\begin{theorem}[The original exponent of a genuine weighted SNC product]
\label{hir:component-product-original-exponent}
\uses{hir:component-product-coordinate-generator,hir:monomial-case,hir:snc-dimension-zero-empty}
\lean{Hironaka.HironakaSNCComponentCoefficients.divExp_componentPowerProduct_eq}
\leanok
For every pair $(W,a)\in A$ and every original point $x\in W$, the
actual exponent of the complete weighted SNC product along $W$ is $a$:
\[
 e_W\left(\prod_i I_{W_i}^{a_i},x\right)=a.
\]
\end{theorem}
\begin{proof}
The preceding theorem produces the actual coordinate monomial generator
near $x$.  It is divisible by $u_W^a$, so its exponent is at least $a$.
To exclude divisibility by $u_W^{a+1}$, keep $u_W=0$ and vary the other
active coordinates to small nonzero values.  Such points of $W$ occur
arbitrarily close to $x$.  At these points the factors belonging to
other active members are units; the remaining coordinate power has
order exactly $a$ because a coordinate has order one and orders of
nonzero holomorphic products add.  The genuine generator lies in the
product ideal, so that ideal cannot have order at least $a+1$ along
$W$ near $x$.
This is the established principal-monomial exponent calculation,
applied to the actual presentation just produced.  Its lower and upper
bounds give equality, also when $a=0$.
In dimension zero, there is no point of an SNC member by
Lemma~\ref{hir:snc-dimension-zero-empty}, so the quantified membership
case cannot occur.  No connectedness of a member is needed for this
weighted-product calculation.
\end{proof}

\begin{theorem}[Produced coefficients belong to the original divisor]
\label{hir:component-factorization-original-coefficients}
\uses{hir:divisor-supported-finite-monomial,hir:component-product-original-exponent}
\dcref{W{\l}odarczyk 2005, Step 2}
\source{wlodarczyk05-hironaka}
\lean{Hironaka.HironakaSNCComponentCoefficients.exists_componentPowerProduct_eq_original_exponents}
\leanok
Let $F$ be a finite list of distinct nonempty connected closed smooth
SNC hypersurfaces, and let $D$ be an effective divisor supported on
$\bigcup_{W\in F}W$.
There is a weighted list $A$ with underlying list $F$ such that
\[
 I_D=\prod_{(W,a)\in A} I_W^a,
 \qquad e_W(I_D,x)=a\quad((W,a)\in A,\ x\in W).
\]
Both the actual factorization and the coefficient identification are
produced from these geometric hypotheses.
\end{theorem}
\begin{proof}
The genuine extraction and residual-unit theorem constructs the finite
weighted list $A$ and the full original ideal product equality.
Its underlying list is exactly $F$, so distinctness and SNC pass to
the underlying weighted list.
Apply the preceding theorem to this actual product at every original
point of every member.  Substitute the proved ideal equality to obtain
the exponent of the original $D$, rather than of an intermediate
residual.  This proves the two asserted conclusions jointly.
An empty boundary gives the empty weighted list and the unit ideal;
there are then no coefficient instances to check.
Zero coefficients are permitted and mean that the original divisor
has exponent zero on that member.
\end{proof}
